\section{Exact operator contraction for distributed PEPS compression}
\label{sec:peps_distributed_operator}
This is the finite operator representation used in
\cite[Theorem~5.2]{OpenAI2026PolynomialPEPS}, September 24, 2026,
\texttt{04-compression.tex}, lines 565--579.
Let $P$ be a finite party set, $G$ a finite set of gate occurrences,
$P_g\subseteq P$ the participant set of $g$, and $r_g\in P_g$ its chosen root.
Let $\Lambda_g$ be any finite branch alphabet, possibly empty, and
$J_k=\{0,\ldots,k-1\}$ for $k\in\mathbb N$.
The physical input and output alphabets $Y_p,X_p$ may be arbitrary sets;
the finite-dimensional density application uses finite ones.
% This chapter proves representation algebra for supplied party-local matrices.
% Identifying an actual chronological circuit with those matrices and the global
% branch/sample sum, and proving uniform polynomial label bounds, remain open.

\begin{definition}[Active gate occurrences]
    \label{def:compression_operator_active_gates}
    \lean{TNLean.PEPS.Approximation.DistributedOperatorContraction.ActiveGate}
    \leanok
    The active gate occurrences are the elements of the finite set $G$. Occurrences are distinguished even when their operations and participant sets coincide.
\end{definition}

\begin{definition}[Tagged star and sample links]
    \label{def:compression_operator_links}
    \lean{TNLean.PEPS.Approximation.DistributedOperatorContraction.Link}
    \leanok
    The link family $\mathcal L$ is the disjoint union over $g\in G$ of the star labels $q\in P_g\setminus\{r_g\}$ and the unordered distinct pairs in $P_g$. Every link retains both its gate occurrence and its star or sample kind.
\end{definition}

\begin{definition}[Gate occurrences at a party]
    \label{def:compression_operator_party_gates}
    \lean{TNLean.PEPS.Approximation.DistributedOperatorContraction.GateAt}
    \leanok
    For $p\in P$, let $G_p=\{g\in G:p\in P_g\}$. These are precisely the gate labels visible at $p$.
\end{definition}

\begin{definition}[Pair-sample positions]
    \label{def:compression_operator_samples}
    \lean{TNLean.PEPS.Approximation.DistributedOperatorContraction.Sample}
    \lean{TNLean.PEPS.Approximation.DistributedOperatorContraction.SampleAt}
    \leanok
    Let $S=\bigsqcup_{g\in G}\{\{p,q\}\subseteq P_g:p\ne q\}$, retaining the gate occurrence in each position. Put $S_p=\{s\in S:p\in s\}$. Positions do not depend on chosen monomial branches.
\end{definition}

\begin{definition}[Star and sample edges]
    \label{def:compression_operator_edges}
    \lean{TNLean.PEPS.Approximation.DistributedOperatorContraction.starLink}
    \lean{TNLean.PEPS.Approximation.DistributedOperatorContraction.sampleLink}
    \leanok
    For $g\in G$ and $q\in P_g\setminus\{r_g\}$, the star edge is $(g,\mathrm{star},q)\in\mathcal L$. For $s=(g,\{p,q\})\in S$, the sample edge is $(g,\mathrm{sample},\{p,q\})\in\mathcal L$.
\end{definition}

\begin{definition}[Oriented link endpoints]
    \label{def:compression_operator_endpoints}
    \lean{TNLean.PEPS.Approximation.DistributedOperatorContraction.src}
    \lean{TNLean.PEPS.Approximation.DistributedOperatorContraction.tgt}
    \lean{TNLean.PEPS.Approximation.DistributedOperatorContraction.endpoints_eq}
    \leanok
    Orient a star from $r_g$ to $q$. Choose an ordered representative $(s^-,s^+)$ of each sample pair and orient its edge accordingly. For every $\ell\in\mathcal L$, write $u_\ell,v_\ell$ for these endpoints. Then $\{u_\ell,v_\ell\}$ is the party endpoint set of $\ell$.
\end{definition}

\begin{definition}[Party incidences]
    \label{def:compression_operator_incidences}
    \lean{TNLean.PEPS.Approximation.DistributedOperatorContraction.LinkAt}
    \lean{TNLean.PEPS.Approximation.DistributedOperatorContraction.starLinkAt}
    \lean{TNLean.PEPS.Approximation.DistributedOperatorContraction.sampleLinkAt}
    \lean{TNLean.PEPS.Approximation.DistributedOperatorContraction.starGateAt}
    \lean{TNLean.PEPS.Approximation.DistributedOperatorContraction.rootGateAt}
    \leanok
    Put $I_p=\{\ell\in\mathcal L:p\in\{u_\ell,v_\ell\}\}$. The star $(g,\mathrm{star},q)$ belongs to $I_p$ when $p=r_g$ or $p=q$; either case gives $g\in G_p$. Every $g$ belongs to $G_{r_g}$. A sample position in $S_p$ gives its sample edge in $I_p$.
\end{definition}

\begin{definition}[Virtual alphabets]
    \label{def:compression_operator_alphabets}
    \lean{TNLean.PEPS.Approximation.DistributedOperatorContraction.alphabet}
    \lean{TNLean.PEPS.Approximation.DistributedOperatorContraction.alphabetFintype}
    \leanok
    Assign the alphabet $D_\ell=\Lambda_g$ to a star edge of $g$ and $D_\ell=J_k$ to a sample edge. These alphabets are finite, including when empty. Private memory and garbage coordinates are absent from $D_\ell$.
\end{definition}

\begin{definition}[Gate-label assignments]
    \label{def:compression_operator_gate_labels}
    \lean{TNLean.PEPS.Approximation.DistributedOperatorContraction.GateLabels}
    \lean{TNLean.PEPS.Approximation.DistributedOperatorContraction.GateLabelsAt}
    \lean{TNLean.PEPS.Approximation.DistributedOperatorContraction.restrictGates}
    \leanok
    A global gate assignment is $\kappa\in\prod_{g\in G}\Lambda_g$, and a local assignment at $p$ is $\xi_p\in\prod_{g\in G_p}\Lambda_g$. The restriction of $\kappa$ to $G_p$ is denoted $\kappa_p$.
\end{definition}

\begin{definition}[Virtual-label assignments]
    \label{def:compression_operator_link_labels}
    \lean{TNLean.PEPS.Approximation.DistributedOperatorContraction.LinkLabels}
    \lean{TNLean.PEPS.Approximation.DistributedOperatorContraction.LinkLabelsAt}
    \lean{TNLean.PEPS.Approximation.DistributedOperatorContraction.restrictLinks}
    \leanok
    A global virtual assignment is $\eta\in\prod_{\ell\in\mathcal L}D_\ell$. Its restriction $\eta_p$ to $I_p$ is a party-local virtual assignment.
\end{definition}

\begin{definition}[Local sample assignments]
    \label{def:compression_operator_sample_labels}
    \lean{TNLean.PEPS.Approximation.DistributedOperatorContraction.restrictSamples}
    \lean{TNLean.PEPS.Approximation.DistributedOperatorContraction.localSamples}
    \leanok
    A global sample assignment is $j\in J_k^S$, whose restriction to $S_p$ is $j_p$. A local virtual assignment $\eta_p$ gives the sample assignment $j(\eta_p)_s=\eta_p(g,\mathrm{sample},s)$ for $s\in S_p$.
\end{definition}

\begin{definition}[Shared virtual labels]
    \label{def:compression_operator_shared_labels}
    \lean{TNLean.PEPS.Approximation.DistributedOperatorContraction.sharedLinkLabels}
    \leanok
    Given $(\kappa,j)$, define $\eta(\kappa,j)_{(g,\mathrm{star},q)}=\kappa_g$ and $\eta(\kappa,j)_{(g,\mathrm{sample},s)}=j_s$. Thus each star carries the complete common gate label and each sample edge carries its position's sample value.
\end{definition}

\begin{definition}[Local star consistency]
    \label{def:compression_operator_consistency}
    \lean{TNLean.PEPS.Approximation.DistributedOperatorContraction.StarConsistent}
    \leanok
    The pair $(\eta_p,\xi_p)$ is consistent if $\eta_p(g,\mathrm{star},q)=\xi_p(g)$ for every star edge incident to $p$. No condition is imposed on sample or physical indices.
\end{definition}

\begin{theorem}[Consistency of shared assignments]
    \label{thm:compression_operator_shared_consistency}
    \lean{TNLean.PEPS.Approximation.DistributedOperatorContraction.sharedLinkLabels_consistent}
    \leanok
    For every global pair $(\kappa,j)$ and party $p$, the assignments $(\eta(\kappa,j)_p,\kappa_p)$ are consistent.
\end{theorem}

\begin{proof}\leanok
    \uses{def:compression_operator_shared_labels,def:compression_operator_consistency}
    At each incident star both labels are $\kappa_g$.
\end{proof}

\begin{definition}[Root-owned scalar coefficients]
    \label{def:compression_operator_root_coefficients}
    \lean{TNLean.PEPS.Approximation.DistributedOperatorContraction.rootCoefficient}
    \leanok
    Let $c_g:\Lambda_g\to\mathbb C$ be arbitrary deterministic coefficients. For a local assignment $\xi_p$, put $C_p(\xi_p)=\prod_{g:r_g=p}c_g(\xi_p(g))$. A gate with one participant has no star edge, but its label and coefficient are still present at its root.
\end{definition}

\begin{definition}[Local operator tensor]
    \label{def:compression_operator_local_tensor}
    \lean{TNLean.PEPS.Approximation.DistributedOperatorContraction.localTensor}
    \leanok
    Let $X_p,Y_p$ be the output and input alphabets, and fix arbitrary rectangular matrices $T_p(\xi_p,j_p):Y_p\to X_p$. Define the local tensor by
    \begin{align}
        A_p(\eta_p;x_p,y_p)
         & =\sum_{\xi_p\in\prod_{g\in G_p}\Lambda_g}
        \mathbf1_{\{(\eta_p,\xi_p)\ {\rm consistent}\}}
        C_p(\xi_p)\,T_p(\xi_p,j(\eta_p))_{x_p,y_p}.
        \label{eq:compression_operator_local_tensor}
    \end{align}
    No positivity, Hermiticity, normalization or bound on the entries of $T_p$ is required.
\end{definition}

\begin{definition}[Globally consistent configurations]
    \label{def:compression_operator_consistent_configs}
    \lean{TNLean.PEPS.Approximation.DistributedOperatorContraction.ConsistentConfig}
    \leanok
    A globally consistent configuration is a pair $(\eta,(\xi_p)_{p\in P})$ such that $(\eta_p,\xi_p)$ is consistent at every party.
\end{definition}

\begin{definition}[Recovery at gate roots and sample links]
    \label{def:compression_operator_recovery}
    \lean{TNLean.PEPS.Approximation.DistributedOperatorContraction.gateLabelsOfConfig}
    \lean{TNLean.PEPS.Approximation.DistributedOperatorContraction.samplesOfConfig}
    \leanok
    For a globally consistent configuration, define $\kappa_g=\xi_{r_g}(g)$ and let $j_s$ be the virtual label on the sample edge of $s$.
\end{definition}

\begin{theorem}[Every participant agrees with the root]
    \label{thm:compression_operator_root_agreement}
    \lean{TNLean.PEPS.Approximation.DistributedOperatorContraction.gateLabel_eq_root}
    \leanok
    In a globally consistent configuration, $\xi_p(g)=\kappa_g$ whenever $g\in G_p$.
\end{theorem}

\begin{proof}\leanok
    \uses{def:compression_operator_consistency,def:compression_operator_recovery}
    If $p=r_g$ this is the definition of $\kappa_g$. Otherwise $(g,\mathrm{star},p)$ is a star edge, and consistency at its two endpoints gives $\xi_p(g)=\eta_{(g,\mathrm{star},p)}=\xi_{r_g}(g)$.
\end{proof}

\begin{theorem}[Recovery of every virtual label]
    \label{thm:compression_operator_link_recovery}
    \lean{TNLean.PEPS.Approximation.DistributedOperatorContraction.sharedLinkLabels_recover}
    \leanok
    If $(\kappa,j)$ is recovered from a globally consistent configuration $(\eta,(\xi_p))$, then $\eta(\kappa,j)=\eta$.
\end{theorem}

\begin{proof}\leanok
    \uses{def:compression_operator_consistency,def:compression_operator_recovery,def:compression_operator_shared_labels}
    A star edge is determined by consistency at its root, and a sample edge is determined by the definition of $j_s$.
\end{proof}

\begin{definition}[Bijection of consistent configurations]
    \label{def:compression_operator_configuration_bijection}
    \lean{TNLean.PEPS.Approximation.DistributedOperatorContraction.consistentConfigEquiv}
    \leanok
    The map $(\kappa,j)\mapsto(\eta(\kappa,j),(\kappa_p)_{p\in P})$ is a bijection from $(\prod_{g\in G}\Lambda_g)\times J_k^S$ onto the globally consistent configurations. Its inverse is recovery at the participating roots and sample links. The two inverse identities follow from root agreement and recovery of every virtual label. Singleton gates and empty alphabets are included.
\end{definition}

\begin{definition}[Reading endpoint configurations]
    \label{def:compression_operator_endpoint_reader}
    \lean{TNLean.PEPS.Approximation.DistributedOperatorContraction.endpointLabels}
    \leanok
    Give each edge independent labels at its two endpoints. At a party $p$, read the first endpoint when $u_\ell=p$ and the second otherwise. This assigns a local virtual label to every $\ell\in I_p$.
\end{definition}

\begin{theorem}[Reading diagonal endpoint labels]
    \label{thm:compression_operator_diagonal_reader}
    \lean{TNLean.PEPS.Approximation.DistributedOperatorContraction.endpointLabels_diagonal}
    \leanok
    If both endpoints of each edge carry the same global label $\eta_\ell$, then the labels read at $p$ are $\eta_p$.
\end{theorem}

\begin{proof}\leanok
    \uses{def:compression_operator_endpoint_reader,def:compression_operator_link_labels}
    Either endpoint choice reads $\eta_\ell$.
\end{proof}

\begin{theorem}[Samples of a shared assignment]
    \label{thm:compression_operator_shared_samples}
    \lean{TNLean.PEPS.Approximation.DistributedOperatorContraction.localSamples_sharedLinkLabels}
    \leanok
    The sample labels read from $\eta(\kappa,j)_p$ are $j_p$.
\end{theorem}

\begin{proof}\leanok
    \uses{def:compression_operator_shared_labels,def:compression_operator_sample_labels}
    By the definition of the shared assignment, the sample edge of $s$ reads $j_s$.
\end{proof}

\begin{theorem}[Each scalar coefficient appears once]
    \label{thm:compression_operator_coefficient_product}
    \lean{TNLean.PEPS.Approximation.DistributedOperatorContraction.prod_rootCoefficient}
    \leanok
    For every global gate assignment $\kappa$, $\prod_{p\in P}C_p(\kappa_p)=\prod_{g\in G}c_g(\kappa_g)$.
\end{theorem}

\begin{proof}\leanok
    \uses{def:compression_operator_root_coefficients}
    Regrouping the gate product by the fibers of $g\mapsto r_g$ gives the equality because every active gate has a unique root.
\end{proof}

\begin{definition}[Endpoint form of the local tensor]
    \label{def:compression_operator_endpoint_tensor}
    \lean{TNLean.PEPS.Approximation.DistributedOperatorContraction.endpointTensor}
    \leanok
    Apply the endpoint reader to a party's independent endpoint configuration and evaluate $A_p$ on the resulting local virtual labels. The physical argument remains the output-input pair $(x_p,y_p)$.
\end{definition}

\begin{definition}[Contracted operator]
    \label{def:compression_operator_operator}
    \lean{TNLean.PEPS.Approximation.DistributedOperatorContraction.contractedOperator}
    \leanok
    Use identity bond matrices on every $D_\ell$, the endpoint form of $A_p$, and physical output-input pairs. The resulting contraction defines the rectangular operator $\sigma:\prod_pY_p\to\prod_pX_p$.
\end{definition}

\begin{theorem}[Reduction to one label per link]
    \label{thm:compression_operator_link_sum}
    \lean{TNLean.PEPS.Approximation.DistributedOperatorContraction.contractedOperator_eq_linkSum}
    \leanok
    The identity-bond contraction satisfies
    \begin{align}
        \sigma(x,y)
         & =\sum_{\eta\in\prod_{\ell\in\mathcal L}D_\ell}
        \prod_{p\in P}A_p(\eta_p;x_p,y_p).
        \label{eq:compression_operator_link_sum}
    \end{align}
\end{theorem}

\begin{proof}\leanok
    \uses{def:compression_operator_operator,thm:compression_operator_diagonal_reader}
    The identity matrix on each edge forces its two endpoint labels to coincide, and reindexing the remaining diagonal assignments by one label per edge makes the endpoint reader yield $\eta_p$.
\end{proof}

\begin{theorem}[Identity-bond factorization across a vertex partition]
    \label{thm:peps_identity_cut_factorization}
    \lean{TNLean.PEPS.DependentBondNetwork.flattening_network_one_eq_mul,
        TNLean.PEPS.DependentBondNetwork.rank_flattening_network_one_le}
    \leanok
    \uses{def:peps_dependent_network,def:peps_group_contraction}
    Let a finite directed multigraph $(V,E)$ have endpoints $t(e),h(e)$,
    finite virtual alphabets $D_e$, and local tensors $A_v$ with physical
    alphabet $P_v$. Let $Z$ be its contraction with identity bond matrices.
    For $S\subseteq V$, let $\partial S$ be the tagged edges with one endpoint
    in $S$ and one in $S^c$, and set $K_S=\prod_{e\in\partial S}D_e$.
    Let $U_S(x,\eta)$ be the partial contraction on $S$, summing the labels
    of all edges internal to $S$ and keeping $x\in\prod_{v\in S}P_v$ and
    $\eta\in K_S$ open. Define $U_{S^c}$ in the same way.
    Let $F_S$ be the matrix reshaping of $Z$, with rows indexed by the
    physical configurations on $S$ and columns by those on $S^c$. Then
    \begin{align}
        F_S & =U_S U_{S^c}^{\mathsf T}.
        \label{eq:peps_identity_cut_factor}
    \end{align}
    If the physical alphabets $P_v$ are finite, then
    \begin{align}
        \operatorname{rank}F_S
         & \le |K_S|=\prod_{e\in\partial S}|D_e|.
        \label{eq:peps_identity_cut_rank}
    \end{align}
    Thus each tagged crossing edge contributes its own alphabet once.
    This is the finite contraction step in
    \cite[Part~II, proof of Theorem~B]{TNLean2026OneDimensional}.
\end{theorem}

\begin{proof}\leanok
    \uses{def:peps_dependent_network,def:peps_group_contraction}
    Each identity bond equates its two endpoint labels. Write $\widetilde A_v$
    for $A_v$ with the same edge label read at every endpoint incidence,
    including both incidences of a loop. Splitting a global edge assignment
    into its two internal parts $\alpha,\beta$ and its crossing part $\eta$
    gives
    \begin{align}
        F_S(x,y)
         & =\sum_{\eta\in K_S}
        \left(\sum_{\alpha}\prod_{v\in S}
        \widetilde A_v((\alpha,\eta)|_v,x_v)\right)
        \left(\sum_{\beta}\prod_{v\in S^c}
        \widetilde A_v((\beta,\eta)|_v,y_v)\right).
        \label{eq:peps_identity_cut_sum}
    \end{align}
    These are precisely the two partial contractions, proving
    (\ref{eq:peps_identity_cut_factor}). Matrix rank is bounded by the
    number of columns of the left factor, namely $|K_S|$, which gives
    (\ref{eq:peps_identity_cut_rank}). The factors pair by the ordinary transpose.
    A loop is internal to one side. Parallel edges retain distinct labels.
    Empty alphabets are allowed: an empty crossing alphabet makes $K_S$
    empty and the matrix zero; with no crossing edges, $K_S$ is a singleton.
\end{proof}

\begin{theorem}[Cut ranks of distributed product-basis columns]
    \label{thm:compression_operator_column_cut_rank}
    \lean{TNLean.PEPS.Approximation.DistributedOperatorContraction.flattening_contractedOperator_column_eq_mul,
        TNLean.PEPS.Approximation.DistributedOperatorContraction.rank_flattening_contractedOperator_column_le}
    \leanok
    \uses{def:compression_operator_operator,def:compression_operator_links,
        def:compression_operator_alphabets,def:peps_ket_column_dependent_fix_input}
    Let $\sigma$ be the distributed operator contraction with tagged link
    family $\mathcal L$, virtual alphabets $D_\ell$, physical output alphabets
    $X_p$, and input alphabets $Y_p$. For any product-basis input
    $z\in\prod_pY_p$, fix that input locally as in
    Definition~\ref{def:peps_ket_column_dependent_fix_input}.
    For a party subset $S$, write $F_{S,z}$ for the coefficient matrix of
    $x\mapsto\sigma(x,z)$ across $S\mid S^c$, and write
    $U_{S,z},U_{S^c,z}$ for the two partial contractions of these actual
    input-fixed local tensors. Then
    \begin{align}
        F_{S,z} & =U_{S,z}U_{S^c,z}^{\mathsf T}.
        \label{eq:compression_column_cut_factor}
    \end{align}
    If every $X_p$ is finite, then
    \begin{align}
        \operatorname{rank}F_{S,z}
         & \le\prod_{\ell\in\partial S}|D_\ell|.
        \label{eq:compression_column_cut_rank}
    \end{align}
    The edge set, alphabets and identity bond matrices are unchanged by
    fixing $z$, as in \cite[Lemma~8.1]{OpenAI2026PolynomialPEPS}.
\end{theorem}

\begin{proof}\leanok
    \uses{thm:peps_identity_cut_factorization,
        def:peps_ket_column_dependent_fix_input,def:compression_operator_operator}
    The input-fixed local tensor is
    $A_{p,z}(\eta,x_p)=A_p(\eta;(x_p,z_p))$.
    Its identity-bond contraction is exactly $\sigma(x,z)$.
    Apply Theorem~\ref{thm:peps_identity_cut_factorization} to these tensors
    on the same tagged star and sample links. The local sums and the coefficient
    at each chosen root stay inside their original tensors. No normalization
    or injectivity condition is required.
\end{proof}

\begin{theorem}[Exact branch and sample evaluation]
    \label{thm:compression_operator_branch_sum}
    \lean{TNLean.PEPS.Approximation.DistributedOperatorContraction.contractedOperator_eq_branchSampleSum}
    \leanok
    The constructed operator has entries
    \begin{align}
        \sigma(x,y)
         & =\sum_{\kappa\in\prod_{g\in G}\Lambda_g}
        \left(\prod_{g\in G}c_g(\kappa_g)\right)
        \sum_{j\in J_k^S}\prod_{p\in P}T_p(\kappa_p,j_p)_{x_p,y_p}.
        \label{eq:compression_operator_branch_sum}
    \end{align}
    The identity permits empty branch alphabets, empty link families, singleton gates and $k=0$.
\end{theorem}

\begin{proof}\leanok
    \uses{thm:compression_operator_link_sum,def:compression_operator_local_tensor,def:compression_operator_configuration_bijection,thm:compression_operator_shared_samples,thm:compression_operator_coefficient_product}
    Expanding the product of the local sums in~(\ref{eq:compression_operator_link_sum}) via~(\ref{eq:compression_operator_local_tensor}) makes the consistency indicators restrict the sum to globally consistent configurations. Reindexing by the bijection with $(\kappa,j)$ and applying the shared-sample identity together with the root-coefficient product identity then gives the display.
\end{proof}

\begin{theorem}[Rectangular matrix form of exact evaluation]
    \label{thm:compression_operator_matrix_sum}
    \lean{TNLean.PEPS.Approximation.DistributedOperatorContraction.contractedOperator_eq_sum_productMatrix}
    \leanok
    The operator identity corresponding to~(\ref{eq:compression_operator_branch_sum}) is
    \begin{align}
        \sigma
         & =\sum_{\kappa\in\prod_{g\in G}\Lambda_g}
        \left(\prod_{g\in G}c_g(\kappa_g)\right)
        \sum_{j\in J_k^S}\bigotimes_{p\in P}T_p(\kappa_p,j_p).
        \label{eq:compression_operator_matrix_sum}
    \end{align}
    The tensor product here has entry $\prod_pT_p(\kappa_p,j_p)_{x_p,y_p}$ and allows different input and output alphabets at each party.
\end{theorem}

\begin{proof}\leanok
    \uses{thm:compression_operator_branch_sum}
    Evaluating both sides at $(x,y)$, distributing the finite matrix sums and scalar multiplication, and applying the entrywise identity gives the operator equality.
\end{proof}

\begin{theorem}[Exact virtual dimensions]
    \label{thm:compression_operator_alphabet_card}
    \lean{TNLean.PEPS.Approximation.DistributedOperatorContraction.card_alphabet}
    \leanok
    A star edge belonging to $g$ has dimension $|\Lambda_g|$, and every sample edge has dimension $k$. These dimensions contain no private-memory or source-support factor.
\end{theorem}

\begin{proof}\leanok
    \uses{def:compression_operator_alphabets}
    The two cardinalities follow directly from the two cases in the definition of $D_\ell$.
\end{proof}

\begin{theorem}[Uniform virtual dimension bound]
    \label{thm:compression_operator_alphabet_bound}
    \lean{TNLean.PEPS.Approximation.DistributedOperatorContraction.card_alphabet_le}
    \leanok
    If $|\Lambda_g|\le M$ for every $g\in G$, then every virtual edge satisfies $|D_\ell|\le\max\{M,k\}$.
\end{theorem}

\begin{proof}\leanok
    \uses{thm:compression_operator_alphabet_card}
    A star edge inherits the given bound at its active gate, and a sample edge has dimension $k$.
\end{proof}

\begin{definition}[Density-branch coefficient]
    \label{def:compression_operator_density_coefficient}
    \lean{TNLean.PEPS.Approximation.DistributedOperatorContraction.densityPairCoefficient}
    \leanok
    Let $\Lambda_g=K_g\times B_g$ be the complete ordered ket/bra label pair. For coefficient functions $a_g,b_g$ and a gate-owned weight $w_g$, put $c_g(\xi,\zeta)=w_g a_g(\xi)\overline{b_g(\zeta)}$. Any averaging factor for a sample position is included once in its gate's $w_g$.
\end{definition}

\begin{theorem}[Density-pair specialization]
    \label{thm:compression_operator_density_pair}
    \lean{TNLean.PEPS.Approximation.DistributedOperatorContraction.contractedOperator_densityPair_eq}
    \leanok
    For arbitrary finite $K_g,B_g$, the star shares the entire pair $(\xi_g,\zeta_g)$. The contracted operator is
    \begin{align}
        \sigma
         & =\sum_{(\xi,\zeta)\in\prod_g(K_g\times B_g)}
        \left(\prod_g w_g a_g(\xi_g)\overline{b_g(\zeta_g)}\right)
        \sum_{j\in J_k^S}\bigotimes_p T_p((\xi,\zeta)_p,j_p).
        \label{eq:compression_operator_density_pair}
    \end{align}
    Each gate-owned weight appears once, including at a singleton gate.
\end{theorem}

\begin{proof}\leanok
    \uses{thm:compression_operator_matrix_sum,def:compression_operator_density_coefficient}
    Specializing the rectangular matrix identity to the complete ket/bra alphabet and the stated coefficient function gives the claim.
\end{proof}

\begin{theorem}[A scalar in one vertex tensor]
    \label{thm:peps_dependent_network_scalar}
    \lean{TNLean.PEPS.DependentBondNetwork.network_mul_at_vertex}
    \leanok
    \uses{def:peps_dependent_network}
    Multiplying the tensor $A_p$ at one vertex $p$ by $s\in\C$ and
    leaving the other tensors unchanged multiplies
    the contraction $\mathcal N_A(B)(\sigma)$ of
    Definition~\ref{def:peps_dependent_network} by $s$. The links and their dimensions
    are unchanged.
\end{theorem}

\begin{proof}
    \leanok
    Every summand of the contraction in
    Definition~\ref{def:peps_dependent_network} contains exactly one
    factor $A_p$.
\end{proof}

\subsection{Finite coordinates for the distributed network}

\begin{definition}[Enumerated virtual alphabets]
    \label{def:compression_operator_finite_coordinates}
    \lean{TNLean.PEPS.Approximation.DistributedOperatorContraction.toPartyNetwork}
    \leanok
    \uses{def:compression_operator_local_tensor,def:peps_party_network_routing}
    In a finite change of coordinates for the construction in
    \cite[proof of Theorem~5.2 and Lemma~8.2]{OpenAI2026PolynomialPEPS}, use the
    distributed links and their finite alphabets $A_e$ from
    Definition~\ref{def:compression_operator_operator}. For each link choose
    a bijection $\vartheta_e:A_e\to[|A_e|]$, where
    $[m]=\{0,\ldots,m-1\}$. Fix $q\in\mathbb N$ and physical readouts
    $u_p:[q]\to X_p$ and $v_p:[q]\to Y_p$ at each party $p$.
    The enumerated network has the same links and endpoints, virtual
    dimension $|A_e|$ on $e$, and local tensors
    \begin{align}
        \widetilde A_p(\eta_p;s,t)
         & = A_p\bigl((\vartheta_e^{-1}(\eta_{p,e}))_{e\ni p};
        u_p(s),v_p(t)\bigr).
        \label{eq:compression_finite_coordinates}
    \end{align}
    Here $(s,t)$ is read as one index in $[q^2]$.
\end{definition}

\begin{theorem}[Contraction in enumerated coordinates]
    \label{thm:compression_operator_finite_coordinates}
    \lean{TNLean.PEPS.Approximation.DistributedOperatorContraction.operatorCoeff_toPartyNetwork}
    \leanok
    \uses{def:compression_operator_finite_coordinates,def:compression_operator_operator}
    Let $V$ be any set, and assign a site $a_p\in V$ to each party.
    For $x,y:V\to[q]$, the operator of the network in
    Definition~\ref{def:compression_operator_finite_coordinates}, read at
    these sites, satisfies
    \begin{align}
        \widetilde\sigma(x,y)
         & = \sigma\bigl((u_p(x(a_p)))_p,(v_p(y(a_p)))_p\bigr).
        \label{eq:compression_coordinate_operator}
    \end{align}
    The readouts and site assignment need not be injective or surjective.
\end{theorem}
\begin{proof}\leanok
    \uses{def:compression_operator_finite_coordinates,thm:compression_operator_link_sum}
    The product bijection $\prod_e\vartheta_e$ changes variables in the
    finite contraction sum. By (\ref{eq:compression_finite_coordinates}),
    each transformed local factor is the corresponding factor in
    (\ref{eq:compression_operator_link_sum}) at the prescribed physical
    coordinates. Their products and hence their sums agree.
\end{proof}
